\begin{tikzpicture}[
  font=\small,
  block/.style={draw, rounded corners=2pt, align=center, minimum height=0.85cm, minimum width=2.5cm, fill=gray!8},
  mask/.style={draw, rounded corners=2pt, align=center, minimum height=0.85cm, minimum width=2.5cm, fill=purple!10, draw=purple!55!black},
  arrow/.style={-{Latex[length=2mm]}, thick}
]
\node[block] (input) {Input frame\\with mixed defects};
\node[mask, right=0.8cm of input] (scratch) {High-frequency response\\scratches and dust};
\node[mask, right=0.8cm of scratch] (flicker) {Low-frequency response\\global flicker};
\node[block, right=0.8cm of flicker] (spatial) {Spatial-domain\\RTN indicator};
\node[block, below=0.8cm of flicker] (fusion) {Confidence-weighted mask fusion};
\node[block, right=0.8cm of fusion] (out) {Guided restoration\\with reduced false alarms};

\draw[arrow] (input) -- (scratch);
\draw[arrow] (scratch) -- (flicker);
\draw[arrow] (flicker) -- (spatial);
\draw[arrow] (scratch) -- (fusion);
\draw[arrow] (flicker) -- (fusion);
\draw[arrow] (spatial) -- (fusion);
\draw[arrow] (fusion) -- (out);

\node[font=\footnotesize, align=center, below=0.25cm of input] {The two wavelet responses separate defect types before fusion.};
\node[font=\footnotesize, align=center, below=0.25cm of out] {The output remains a degradation indicator, not a hallucinated image prior.};
\end{tikzpicture}
